% !PRINTED 320
\weberpage{320}
Wendet man auf die Indices $\xi, \eta$ der $X$ die Substitution $\sigma$
an, so erleiden die Indices der $\Psi$ [wie im \S.\,80, (2)] die
Substitution $\varrho^{-1}$, wenn $\varrho$ die transponirte Substitution
von $\sigma$ ist.

Es ist aber, den Substitutionen (4) entsprechend,
\weq{6}{\varrho_1^{-1} = \begin{pmatrix}0, & 1\\1, & 0\end{pmatrix}, \quad
\varrho_2^{-1} = \begin{pmatrix}1, & 0\\1, & 1\end{pmatrix}.}

Indem wir nun der Einfachheit halber voraussetzen, dass die drei Grössen
(5) von Null verschieden sind, setzen wir
\weq{7}{\begin{aligned}
v_{1,0} &= \frac{\Psi_{0,1}\Psi_{1,1}}{\Psi_{1,0}}
= \frac{\Psi_{0,1}\Psi_{1,0}\Psi_{1,1}}{\Psi_{1,0}^2};\\
v_{0,1} &= \frac{\Psi_{1,0}\Psi_{1,1}}{\Psi_{0,1}}
= \frac{\Psi_{0,1}\Psi_{1,0}\Psi_{1,1}}{\Psi_{0,1}^2};\\
v_{1,1} &= \frac{\Psi_{1,0}\Psi_{0,1}}{\Psi_{1,1}}
= \frac{\Psi_{1,0}\Psi_{0,1}\Psi_{1,1}}{\Psi_{1,1}^2}.
\end{aligned}}

Die Grössen $v_{\xi,\eta}$ bleiben durch die Substitutionen von $Q$
ungeändert und erleiden durch $\sigma_1, \sigma_2$ die durch (6)
bestimmten Permutationen
\[
\begin{array}{l l l l}
 & v_{1,0}, & v_{0,1}, & v_{1,1}\\
\sigma_1 & v_{0,1}, & v_{1,0}, & v_{1,1}\\
\sigma_2 & v_{1,1}, & v_{0,1}, & v_{1,0}.
\end{array}
\]

Daraus ergiebt sich, dass die symmetrischen Functionen der drei Grössen
$v_{1,0}, v_{0,1}, v_{1,1}$ durch die Permutationen $P$ ungeändert
bleiben, und folglich in dem Rationalitätsbereiche der symmetrischen
Functionen der $X$ enthalten sind. Die drei Grössen $v$ sind also die
Wurzeln einer cubischen Gleichung, und es würde keine Schwierigkeit
haben, diese cubische Resolvente der allgemeinen biquadratischen
Gleichung zu bilden.

Aus (7) folgt aber ferner
\[
v_{1,0}\,v_{0,1} = \psi_{1,1}^2, \quad v_{1,0}\,v_{1,1} = \psi_{0,1}^2,
\quad v_{0,1}\,v_{1,1} = \psi_{1,0}^2, % [sic] printed with lowercase psi, elsewhere capital Psi is used for these quantities
\]
und demnach ergiebt sich aus (5), wenn wir noch die rationale Grösse
\[
X_{0,0} + X_{1,0} + X_{0,1} + X_{1,1} = a
\]
setzen,
\weq{8}{4X_{0,0} = a + \sqrt{v_{1,0}}\,\sqrt{v_{0,1}}
+ \sqrt{v_{1,0}}\,\sqrt{v_{1,1}} + \sqrt{v_{0,1}}\,\sqrt{v_{1,1}}.}

Der Ausdruck (8) giebt nur vier verschiedene Werthe, wenn wir die
Vorzeichen der drei darin vorkommenden Quadratwurzeln auf alle mögliche
Arten bestimmen, und man kann auch um-
%CONT>>
